%----------------------------------------------------------------------------------------
%	PACKAGES AND OTHER DOCUMENT CONFIGURATIONS
%----------------------------------------------------------------------------------------

\documentclass[10pt,stdletter,dateno,sigleft]{newlfm} % Extra options: 'sigleft' for a left-aligned signature, 'stdletternofrom' to remove the from address, 'letterpaper' for US letter paper - consult the newlfm class manual for more options

%\usepackage{charter} % Use the Charter font for the document text
\usepackage{graphics}

\usepackage{standalone}

%% with this patch, the newlfm should work (BEGIN)
\usepackage{etoolbox}
\makeatletter
\patchcmd{\@zfancyhead}{\fancy@reset}{\f@nch@reset}{}{}
\patchcmd{\@set@em@up}{\f@ncyolh}{\f@nch@olh}{}{}
\patchcmd{\@set@em@up}{\f@ncyolh}{\f@nch@olh}{}{}
\patchcmd{\@set@em@up}{\f@ncyorh}{\f@nch@orh}{}{}
\makeatother
%% with this patch, the newlfm should work (END)


\newsavebox{\Luiuc}\sbox{\Luiuc}{\parbox[b]{1.75in}{\vspace{0.5in}
\includegraphics[width=1.2\linewidth]{/Users/hectorbahamonde/Job_Market/logo.png}}} % Company/institution logo at the top left of the page
\makeletterhead{Uiuc}{\Lheader{\usebox{\Luiuc}}}

\newlfmP{sigsize=50pt} % Slightly decrease the height of the signature field
\newlfmP{addrfromphone} % Print a phone number under the sender's address
\newlfmP{addrfromemail} % Print an email address under the sender's address
\PhrPhone{\texttt{p}} % Customize the "Telephone" text
\PhrEmail{\texttt{e}} % Customize the "E-mail" text

\lthUiuc % Print the company/institution logo

% email input begin
\newread\fid
\newcommand{\readfile}[1]% #1 = filename
{\bgroup
  \endlinechar=-1
  \openin\fid=#1
  \read\fid to\filetext
  \loop\ifx\empty\filetext\relax% skip over comments
    \read\fid to\filetext
  \repeat
  \closein\fid
  \global\let\filetext=\filetext
\egroup}
\readfile{/Users/hectorbahamonde/Bibliografia_PoliSci/email.txt}
% email input end

\usepackage{color} % for colors
\usepackage{hyperref} % for links

\hypersetup{
  colorlinks = true,
  urlcolor = blue,
  pdfpagemode = UseNone
}

%----------------------------------------------------------------------------------------
%	YOUR NAME AND CONTACT INFORMATION
%----------------------------------------------------------------------------------------

%\namefrom{\vspace{-3cm}Hector Bahamonde, PhD} % Name

\addrfrom{
{\vspace{-2cm}}\\ % Date
[12pt]
{\bf Hector Bahamonde, PhD}\\
{\color{blue}{\bf Assistant Professor}}\\
\emph{{\color{blue}Social Sciences Institute}}\\
{\color{blue}O$'$Higgins University}\\
Rancagua, Chile
}

\phonefrom{+1(504) 941-9131} % Phone number

\emailfrom{\href{mailto:\filetext}{\filetext} \\ 
\texttt{w}: \href{http://www.hectorbahamonde.com}{www.HectorBahamonde.com}\\
\today, download latest version \href{http://github.com/hbahamonde/Job_Market/raw/master/Bahamonde_Cover_Letter.pdf}{\texttt{{\color{red}here}}}}% Link to last version


%----------------------------------------------------------------------------------------
%	ADDRESSEE AND GREETING/CLOSING
%----------------------------------------------------------------------------------------

%\greetto{\vspace{-1cm}Dear Members of the Search Committee,} % Greeting text
%\closeline{\vspace{-1cm}\includegraphics[width=3cm]{/Users/hectorbahamonde/Administracion/signature.pdf} \hspace{-4cm}} % Closing text




%----------------------------------------------------------------------------------------

\begin{document}
\begin{newlfm}

%----------------------------------------------------------------------------------------
%	LETTER CONTENT
%----------------------------------------------------------------------------------------



\vspace{-2cm}{\bf \huge Research Statement}\\

%\vspace{-1cm}
Within comparative politics, my research explores the intersection between inequality and political development, from a political economy perspective. %My {\bf job market paper} (\emph{\input{/Users/hectorbahamonde/Dissertation/Papers/Earthquake_Paper/status.txt}\unskip}) is entitled ``\href{https://github.com/hbahamonde/Earthquake_Paper/raw/master/Bahamonde_Earthquake_Paper.pdf}{\input{/Users/hectorbahamonde/Dissertation/Papers/Earthquake_Paper/title.txt}\unskip}.'' \input{/Users/hectorbahamonde/Dissertation/Papers/Earthquake_Paper/abstract.txt}\unskip


% Vote Buying in Brazil// explaining the paper
In ``\href{https://journals.sub.uni-hamburg.de/giga/jpla/article/view/1121/1128}{Aiming Right at You: Group versus Individual Clientelistic Targeting in Brazil}'' (Journal of Politics in Latin America) I used {\bf matching techniques} for observational data. \input{/Users/hectorbahamonde/research/Clientelism_paper/abstract.txt}\unskip 


In addition to that, and with the support of a generous grant, I designed in {\bf Qualtrics} two {\bf experiments} in the U.S. out of a series of experiments to be fielded in Latin America for further comparison. In ``\href{https://github.com/hbahamonde/Vote_Selling/raw/master/Bahamonde_VoteSellingUS.pdf}{Still for Sale: The Micro-Dynamics of Vote Selling in the United States, Evidence From a List Experiment}'' (Acta Politica), I looked at the tipping points at which a sample of U.S. citizens (N = 1,479) prefer a monetary incentive rather than keeping their right to choose whom to vote for. \input{/Users/hectorbahamonde/research/Vote_Selling/abstract.txt}\unskip 

The second experiment is a {\bf conjoint experiment} and it is being analyzed in a separate paper entitled ``\href{https://github.com/hbahamonde/Conjoint_US/raw/master/Bahamonde_Quininao_Conjoint.pdf}{Vote Selling in the United States: Introducing Support Vector Machine Methods to Analyzing Conjoint Experimental Data}'' (\emph{\input{/Users/hectorbahamonde/research/Conjoint_US/status.txt}\unskip}). My coauthor and myself are developing a new methodology to analyzing conjoint experiments via {\bf machine learning} techniques, particularly, {\bf support vector machines}. {\input{/Users/hectorbahamonde/research/Conjoint_US/abstract.txt}\unskip}

Finally, and within the same project, we are currently designing an {\bf economic experiment} in the lab. \input{/Users/hectorbahamonde/research/Economic_Experiment_Vote_Selling/abstract.txt}\unskip



% main argument of the book
{\bf My book manuscript} entitled \emph{``Structural Transformations in Latin America: State Building and Elite Competition 1850-2010''} explains that the economic structural transformation in Latin America---the secular decline of agriculture and substantial expansion of manufacturing---imposed tight constraints on the way politics was run by the incumbent landowning class. This was a major change due to the advantage the landed elites enjoyed since colonial times. Where the expansion of the industrial sector was weak, post-colonial norms persisted due to institutional inertia, perpetuating the advantaged position of the agricultural class. Leveraging economic sectoral outputs dating back to 1900 until 2010, for a sample of Latin American countries, I use {\bf panel data methods} (particularly, Cox-proportional hazard models), time series analyses (VAR models, impulse response functions, and Granger-causality tests), and fine-grained qualitative data to support my argument. {\bf Another major contribution of the book} is in measurement. One of the biggest gaps in the literature is the lack of a measurement of state capacities able to capture variations of \emph{stateness} over time. Using a novel dataset, the book proposes measuring state capacities using earthquake data. The rationale is very intuitive: the capacity to enforce quake-sensitive building codes throughout the territory is a \emph{reflection} of the overall (\emph{in})capacity of states of solving both logistic and political limitations at the subnational level. \emph{Why does a 7.0 magnitude earthquake flatten Haiti. leaving at least 100,000 deaths, while a 8.8 earthquake in Chile in the same year leaves just 525 deaths?} By exploiting  variations on earthquake death-tolls and local population to weight the number of deaths since 1900, I measure state capacities over time. Keeping magnitudes constant, casualty differentials should be attributed to the \emph{lack of state capacities}. The paper version of the book is entitled ``\href{https://github.com/hbahamonde/Earthquake_Paper/raw/master/Bahamonde_Earthquake_Paper.pdf}{{\input{/Users/hectorbahamonde/Dissertation/Papers/Earthquake_Paper/title.txt}\unskip}}'' (\emph{\input{/Users/hectorbahamonde/Dissertation/Papers/Earthquake_Paper/status.txt}\unskip}).





In one of the {\bf sections of the book}, I explain that the structural transformation required both sectors to grow in a \emph{balanced} fashion, \emph{leveling both elites in their relative political, economic, and military capacities}. Leveraging the dual sector model of economic growth, I sketch a theory of political and economic development that stresses the structural economic dependence of both sectors. The agricultural sector supplied labor and cheap foodstuff---which the industrial sector demanded, promoting balanced economic development of both sectors. The political consequence of balanced economic growth was the mutual political dependence and the need of inter-elite compromises that fostered both economic and political development in the long-run. The paper version of the chapter (currently \emph{\input{/Users/hectorbahamonde/Dissertation/Papers/NegativeLink/status.txt}\unskip}) can be downloaded \href{https://github.com/hbahamonde/Negative_Link_Paper/raw/master/Bahamonde_NegativeLink.pdf}{\texttt{here}}. 


{\bf Summary and Future Research} In summary, my book manuscript, job market paper and working papers on vote buying and experiments on vote selling work toward exploring the effects of inequality in political development from a comparative perspective. I use a broad methodological approach, historical comparisons, time series analyses and experimental and quasi-experimental methodologies. My goal is to use this toolkit to keep asking ``big'' questions that are fundamental to our discipline. 

My future research will seek to study the connection between elite competition and democratic regimes, exploring the connection between state building and democratic institutions. The first paper within this wave of publications is entitled ``\href{https://github.com/hbahamonde/Inequality_State_Capacity/raw/main/Bahamonde_Trasberg.pdf}{\input{/Users/hectorbahamonde/research/Inequality_State_Capacity/title.txt}\unskip}'' ({\input{/Users/hectorbahamonde/research/Inequality_State_Capacity/status.txt}\unskip}). This study argues that democratic rule combined with high state capacity pro-duce higher levels of income inequality over time. This relationship operates through the positive effect of high-capacity democratic context on foreign direct investment and financial development. The second paper is entitled ``\href{https://raw.githubusercontent.com/hbahamonde/Tobalaba/main/abstract.txt}{\input{/Users/hectorbahamonde/research/Tobalaba/title.txt}\unskip}'' ({\input{/Users/hectorbahamonde/research/Tobalaba/status.txt}\unskip}). {\input{/Users/hectorbahamonde/research/Tobalaba/abstract.txt}\unskip} 

Please check my \href{http://www.hectorbahamonde.com/teaching/}{\texttt{teaching portfolio}} and see how my {\bf research and teaching interests match}. {More information, \href{http://www.hectorbahamonde.com/teaching/}{\texttt{syllabi}}, my \href{http://github.com/hbahamonde/Job_Market/raw/master/Bahamonde_Research_Statement.pdf}{\texttt{research}}, \href{http://github.com/hbahamonde/Job_Market/raw/master/Bahamonde_Teaching_Statement.pdf}{\texttt{teaching}} and \href{http://github.com/hbahamonde/Job_Market/raw/master/Bahamonde_Diversity_Statement.pdf}{\texttt{diversity}} statements, as well as other \href{http://www.hectorbahamonde.com/research/}{\texttt{papers}} are available on my website: \href{http://www.hectorbahamonde.com}{\texttt{www.HectorBahamonde.com}}. Thank you for considering my application. I look forward to hearing from you.\unskip}


\vspace{-10cm}


%----------------------------------------------------------------------------------------

\end{newlfm}
\end{document}
